\section{Methods}
\label{sec:methods}

This section defines how the co-detections are associated, how the
dispersion-measure and scattering budgets are modeled, and how the fitted
quantities used in Section~\ref{sec:results} are validated.

\subsection{Time-of-Arrival Crossmatching}
\label{sec:toa}

Before attributing the dispersion and scattering budgets we verify that the
twelve events are genuine co-detections---the same burst seen by both
facilities, with arrival times consistent under the cold-plasma dispersion law
and the geometric light-travel delay between the sites. The wide baseline
between CHIME/FRB (400--800\,MHz) and DSA-110 ($\sim$1.4\,GHz) makes the
co-detections a test of the $1/\nu^2$ dispersion scaling and of the relative
timing alignment of the two facilities \citep{Lorimer2007, Thornton2013}.

\subsubsection{Delay model}

We translate the arrival times at the two central frequencies to a common
reference frequency $\nu_{\mathrm{ref}}=400$\,MHz. The cold-plasma dispersion
delay is
\begin{equation}
\Delta t_p = K_{\mathrm{DM}}\,\mathrm{DM}
\left(\frac{1}{\nu_{\mathrm{ref}}^2}-\frac{1}{\nu_{\mathrm{obs}}^2}\right),
\label{eq:dmdelay}
\end{equation}
with $K_{\mathrm{DM}} = 4.148808\times10^{3}\,\mathrm{MHz^2\,pc^{-1}\,cm^3\,s}$.
The geometric light-travel difference between the observatories is
\begin{equation}
\tau_{\mathrm{geo}} = \frac{(\mathbf{p}_2-\mathbf{p}_1)\cdot\hat{\mathbf{s}}}{c},
\label{eq:geodelay}
\end{equation}
where $\mathbf{p}_1$ and $\mathbf{p}_2$ are the GCRS position vectors of CHIME
and DSA-110 and $\hat{\mathbf{s}}$ is the unit vector toward the source. The
timing residual is the difference between the frequency-corrected inter-site
offset and the predicted geometric delay; a residual of zero indicates perfect
agreement.

\subsubsection{Residuals and systematics}

Table~\ref{tab:sample} restores the association diagnostics after the V6
re-validation pass. The tabulated residuals use the shared DSA-DM convention:
the CHIME and DSA arrivals are both referenced to 400\,MHz with the DSA catalog
DM, so the residual is a timing/geometry consistency check rather than an
independent per-telescope-DM comparison. Re-referencing the CHIME arrivals with
their independently measured DMs would inject millisecond-scale shifts for the
largest CHIME--DSA DM offsets, so those shifts are kept in the DM-provenance
audit rather than folded into the association residual.

All twelve bursts have chance-coincidence probabilities
$P_{\rm cc}<10^{-8}$ under the conservative baseline association window and
pass the positional-coincidence check. The eight bursts with constrained
CHIME-side DMs also pass the CHIME--DSA DM-agreement check; the four remaining
bursts are marked as position-only associations because their CHIME sub-band
fits did not provide a constraining independent DM.

% Removed figure candidate:
% - Artifact: figures/toa_crossmatch_analysis_premium.pdf
% - Generator: pipeline/crossmatching/plotting.py, plot_toa_analysis()
% - Inputs: pipeline/crossmatching/toa_crossmatch_results.json
% - Reason withheld: all figures are withheld pending intentional re-addition
%   with current trust/provenance comments.
%
% Removed figure candidate:
% - Artifact: figures/systematics_check_matrix.pdf
% - Generator: pipeline/crossmatching/plotting.py, plot_systematics_matrix()
% - Reason withheld: not needed for the scaffolded manuscript. Restore only if a
%   later Results/Discussion pass needs a compact timing-systematics diagnostic.

The independent DM comparison and positional agreement are documented in the V6
validation artifacts; Table~\ref{tab:sample} carries the manuscript-facing
timing residual, $P_{\rm cc}$, and association verdict columns.

\section{The Dispersion and Scattering Budget}
\label{sec:budget}

% Core methods section. The observed DM decomposes as
%   DM_obs = DM_MW,ISM + DM_MW,halo + <DM_cosmic>(z) + DM_intervening + DM_host,
% with DM_MW,ISM from NE2001/YMW16 (pygedm), DM_MW,halo a 40 pc/cm^3 prior
% (Yamasaki & Totani 2020), <DM_cosmic>(z) the Macquart relation
% (Macquart et al. 2020; f_IGM=0.84, chi_e=0.875), and DM_intervening a two-phase
% (hot mNFW + clumpy cool) CGM column for each intersecting foreground galaxy,
% reported raw and capped at b = 0.1 R_vir (the galaxy-interior regime).
%
% Scattering: measured tau_1GHz from nested-sampling fits (quality-gated on
% reduced chi^2, with R^2/normality informational only), compared against the
% predicted intervening tau from a thin-screen two-phase model.

\subsection{Dispersion Measure Decomposition}
\label{sec:dm}

\subsection{Scattering Attribution}
\label{sec:scattering}


\subsection{Band-restricted burst energies}
\label{sec:methods-energies}

For each sightline with a spectroscopic host redshift and a well-constrained
per-band amplitude fit, we estimate the isotropic-equivalent burst energy directly
from the joint CHIME--DSA fit, without extrapolating either band's spectrum
beyond where it is constrained.
The fit returns a per-band spectral amplitude and index,
$F_X(\nu) = c_{0,X}\,(\nu/\nu_{\mathrm{ref},X})^{\gamma_X}$, which we place on
an absolute scale with each telescope's per-channel radiometer conversion
$\sigma_{S,X}(\nu) = \mathrm{SEFD}_X(\nu)\,/\,[\sqrt{n_{\mathrm{pol}}\,\Delta\nu\,\Delta t}\,G_X(\nu)]$,
taking the system-equivalent flux density and primary-beam gain $G_X$ from the
DSA-110 \citep{Law2024} and CHIME/FRB \citep{CHIMEFRB2018,Michilli2021} flux
scales. We integrate over each instrument's observing band---CHIME over
$0.400$--$0.800\,\mathrm{GHz}$ and DSA over $1.311$--$1.499\,\mathrm{GHz}$,
\begin{equation}
E_{\mathrm{iso}} = \frac{4\pi D_L^2(z)}{1+z}\left[
  \int_{\nu_1^{C}}^{\nu_2^{C}} s_C\,F_{\mathrm{CHIME}}(\nu)\,d\nu
  + \int_{\nu_1^{D}}^{\nu_2^{D}} s_D\,F_{\mathrm{DSA}}(\nu)\,d\nu
\right],
\label{eq:eiso}
\end{equation}
with $D_L(z)$ from the fiducial \textit{Planck}\,2018 cosmology
\citep{Planck2018}, the $(1+z)$ bandwidth k-correction applied, and
$1\,\mathrm{Jy\,ms\,Hz} = 10^{-29}\,\mathrm{J\,m^{-2}}$. This band-restricted
integral constrains the energy released within the observed spectral envelope:
it avoids the large, model-dependent extrapolation that a single power-law fit
across the full $0.4$--$1.5\,\mathrm{GHz}$ span would require, at the cost of not
counting flux in the unobserved $0.8$--$1.3\,\mathrm{GHz}$ gap.

Two properties of this definition must be stated prominently. First, the
observed bands map to \emph{per-burst} rest-frame intervals: CHIME's
$0.400$--$0.800\,\mathrm{GHz}$ samples the rest-frame band
$[0.400(1+z),\,0.800(1+z)]\,\mathrm{GHz}$ and DSA's
$1.311$--$1.499\,\mathrm{GHz}$ samples $[1.311(1+z),\,1.499(1+z)]\,\mathrm{GHz}$,
so the rest-frame window shifts from burst to burst. These band-restricted
energies are therefore \emph{not} comparable to literature energies computed
over a fixed rest-frame band, and cannot be placed on a standard FRB luminosity
function without a spectral-model extrapolation that we deliberately avoid here;
we report them as observed-band quantities only. Second, the single $(1+z)$ in
Eq.~\ref{eq:eiso} is the bandwidth (time-dilation) k-correction, which is common
to both bands; the differing per-band spectral indices $\gamma_C,\gamma_D$ enter
only the in-band flux integrals $\int F_X(\nu)\,d\nu$, not the k-correction, so
the two-band sum is internally consistent despite the single $(1+z)$ factor.
% TODO(energies-fixed-band-variant): Optional under V3 — if luminosity-function
% comparison is needed, elect a fixed rest-frame-band companion E_iso (define
% the band, extrapolation across the 0.8--1.3 GHz gap, and table column), then
% restore a Methods sentence. D5 lock (2026-07-10) dropped the unconditional
% promise; see docs/rse/specs/decision-d2-d5-scattering-design-locks.md.
Energy rows are deferred until the spectral
amplitudes, calibration path, host-redshift provenance, and inclusion rule are
ready to be described and cited in the manuscript.
